\chapter{2006:Terence Tao}\label{ch2006terencetao:terence-tao-fields-medal-2006-a-comprehensive-report-on-his-work-and-impact}

\begin{center}\rule{0.5\linewidth}{0.5pt}\end{center}

\textbf{Prepared:} July 2026 \textbf{Subject:} Terence Chi-Shen Tao (b. July 17, 1975, Adelaide, Australia) \textbf{Position:} Professor of Mathematics, University of California, Los Angeles (since 1999) \textbf{Doctoral Advisor:} Elias M. Stein (Princeton University) \textbf{Fields Medal Citation:} \emph{``For his contributions to partial differential equations, combinatorics, harmonic analysis and additive number theory.''}

\begin{center}\rule{0.5\linewidth}{0.5pt}\end{center}

\begin{prereqs} Real analysis ($L^p$ spaces, Fourier analysis, Sobolev spaces), basic PDE theory (well-posedness, energy methods), additive combinatorics (sumset structure, Gowers norms). \end{prereqs}

\begin{objectives} Prove the Green--Tao theorem that primes contain arbitrarily long arithmetic progressions; explain the transference principle; describe the inverse theorem for Gowers norms; understand the Four Moment Theorem for random matrix universality. \end{objectives}

\section{Problem History}

The question of arithmetic progressions in the primes dates to the nineteenth century. Legendre conjectured that there are infinitely many prime pairs in arithmetic progression, and Dirichlet's theorem on primes in arithmetic progressions (1837) confirmed this for progressions with coprime first term and common difference. By the mid-twentieth century, it was widely believed that the primes contain arbitrarily long arithmetic progressions, but no proof was in sight because the primes are a sparse set with density tending to zero, making them inaccessible to density-based methods.

The major obstacle was Szemer\'edi's theorem (1975): any subset of the integers with positive upper density contains arbitrarily long arithmetic progressions. While revolutionary, Szemer\'edi's theorem is fundamentally a statement about dense sets, and the primes have zero upper density. Extending such results to sparse sets requires fundamentally new ideas.

In harmonic analysis, the Kakeya conjecture---that sets in $\mathbb{R}^n$ containing a unit line segment in every direction must have Hausdorff dimension $n$---had resisted proof for decades. The restriction problem, concerning when the Fourier transform of a function on $\mathbb{R}^n$ can be meaningfully restricted to a lower-dimensional surface, was closely related but equally difficult.

In random matrix theory, the Wigner--Dyson--Mehta universality conjecture---that the local eigenvalue spacing distribution of random matrices is independent of the entry distribution---was known only for Gaussian entries. Proving universality for general entries required new techniques.

The \emph{Hardy--Littlewood circle method}, a powerful analytic technique in additive number theory dating to the 1920s, provides another route to understanding additive patterns in the primes. The method decomposes the integers using exponential sums and estimates the major and minor arcs separately. While effective for representations of integers as sums of a bounded number of powers, the circle method encounters fundamental difficulties with the primes because the characteristic function of the primes has poor Fourier analytic properties. The \emph{Vinogradov mean value theorem} (proved by Bourgain, Demeter, and Guth in 2015) provides sharp estimates for the number of solutions to systems of diagonal equations, and it is closely related to Tao's work on the restriction problem and decoupling.

The \emph{restriction problem} asks for the boundedness of the Fourier transform restricted to a curved hypersurface in $\mathbb{R}^n$. The Tomas--Stein restriction theorem (1975) gives the sharp bound for the $L^p \to L^q$ restriction operator, and Tao's work with Guth, Katz, and others extended these results to more general surfaces. The restriction conjecture, which predicts the sharp exponent for the restriction operator, remains open in general, but Tao's methods have produced significant partial results. The restriction problem is closely connected to the Kakeya conjecture, the wave packet decomposition, and the multilinear restriction estimate.

The \emph{Kakeya conjecture} asserts that sets in $\mathbb{R}^n$ containing a unit line segment in every direction must have Hausdorff dimension $n$. The conjecture is known in two dimensions (due to Davies, 1971) but remains open for $n \ge 3$. Tao's work with Guth on the polynomial method produced significant progress on the Kakeya problem, showing that the Minkowski dimension of a Kakeya set in $\mathbb{R}^3$ is at least $5/2 + \varepsilon$ for some $\varepsilon > 0$. This result uses the polynomial partitioning technique of Guth and Katz, which is a powerful tool for studying incidence geometry and combinatorial geometry.

The \emph{Szemer\'edi regularity lemma} (1975) and the \emph{graph removal lemma} provide combinatorial tools for detecting structure in dense graphs and hypergraphs. Tao's work connects these combinatorial tools to additive number theory through the transference principle, showing that pseudorandom sets behave like dense sets for the purposes of detecting additive patterns. The key insight is that the ``error terms'' in the regularity lemma can be controlled uniformly for pseudorandom measures, allowing the density-based conclusions to be transferred. The transference principle has been refined by Green and Tao into a general framework: given a sparse set $A$ and a dense ``model'' $B$, the distribution of linear patterns in $A$ is close to that in $B$, provided $A$ satisfies the appropriate pseudorandomness conditions. The pseudorandomness is measured using the Gowers norm, and the model is constructed using the W-trick and the Hardy--Littlewood circle method.

The \emph{Green--Tao theorem} for polynomial sequences (2010) extends the original result to show that the primes contain arbitrarily long polynomial progressions $p, p+q(1), \ldots, p+q(k)$ where $q(1), \ldots, q(k)$ are integer-valued polynomials. The proof uses a generalization of the transference principle to polynomial sequences, combined with the inverse theorem for Gowers norms and the theory of nilsequences. The key new ingredient is the \emph{linear forms condition}, which ensures that the weighted primes are pseudorandom with respect to all linear forms arising from the polynomial progression.

The \emph{Gowers norms} are defined recursively: $\|f\|_{U^1} = |\mathbb{E}_x f(x)|$ and $\|f\|_{U^k}^{2^k} = \mathbb{E}_h \|f \cdot T_h f\|_{U^{k-1}}^{2^{k-1}}$, where $T_h f(x) = f(x+h)$ is the translation operator. The Gowers $U^2$ norm controls the distribution of 2-term progressions (i.e., the correlation $f(x) f(x+h)$), and the $U^k$ norm controls the distribution of $k$-term progressions. The inverse theorem states that if $\|f\|_{U^k} \ge \delta$, then $f$ has nontrivial correlation with a $(k-1)$-step nilsequence, which is a function of the form $F(g(n)\Gamma)$ where $G$ is a $(k-1)$-step nilpotent Lie group. This inverse theorem is the key to the Green--Tao proof, as it shows that any set with large Gowers norm must contain additive structure.

\section{How Was It Solved}

\textbf{The Green--Tao theorem (2004).} With Ben Green, Tao proved that the primes contain arbitrarily long arithmetic progressions. The proof proceeds in three steps. First, they developed the \emph{transference principle}: a method for applying density-based results (like Szemer\'edi's theorem) to sparse sets that are ``pseudorandom.'' The idea is to find a dense set $B$ that ``models'' the sparse set $A$ in the sense that linear patterns are distributed similarly, apply the density result to $B$, and transfer the conclusion back to $A$.

Second, they proved the pseudorandomness of the primes. Using the W-trick to handle the irregularity of primes (e.g., all large primes are odd), they showed that the weighted primes satisfy the necessary linear forms conditions, verified using classical results on primes in arithmetic progressions. Precisely, the W-trick replaces the primes with the set $\mathcal{P}_w = \{p \in \mathbb{P} : p \equiv w \pmod{W}\}$ where $W = \prod_{p \le N^{1/c}} p$ for some constant $c$, making the set ``uniformly distributed'' in arithmetic progressions.

Third, they applied Szemer\'edi's theorem via the transference principle to conclude that the primes contain arbitrarily long progressions. The proof, over 50 pages, is a landmark in additive number theory. The transference principle has since been refined and extended by many authors, and it now applies to a wide range of sparse sets including polynomial sequences, values of multiplicative functions, and subsets of $\mathbb{F}_q[t]$.

The key technical innovation in the Green--Tao proof is the \emph{W-trick}: replacing the primes $\mathbb{P}$ with the set $\mathcal{P}_w = \{p \in \mathbb{P} : p \equiv w \pmod{W}\}$ for a carefully chosen modulus $W$. The modulus $W$ is chosen to be the product of all primes up to $N^{1/c}$ for a small constant $c$, which makes $\mathcal{P}_w$ ``uniformly distributed'' in arithmetic progressions modulo any fixed $q$. This uniformity is exactly what the transference principle requires. The W-trick introduces a averaging over residue classes $w$ modulo $W$, and the final result is obtained by summing over all admissible $w$.

\textbf{The inverse theorem for Gowers norms.} With Green and Ziegler, Tao proved the inverse conjecture: a function with large Gowers $U^k$ norm $\|f\|_{U^k} \ge \delta$ must have nontrivial correlation with a $(k-1)$-step nilsequence $F(g(n)\Gamma)$, where $G/\Gamma$ is a nilmanifold. The proof, over 200 pages, uses nilmanifold geometry, cohomological methods, and refinements of the Gowers--Cauchy--Schwarz inequality. The Gowers norms are defined recursively: $\|f\|_{U^1} = |\mathbb{E}_x f(x)|$ and $\|f\|_{U^k}^{2^k} = \mathbb{E}_h \|f \cdot T_h f\|_{U^{k-1}}^{2^{k-1}}$.

The inverse theorem for $U^2$ states that if $\|f\|_{U^2} \ge \delta$, then $f$ correlates with a quadratic phase $e^{2\pi i P(n)}$ where $P$ is a polynomial of degree at most 2. For $U^3$, the statement involves 2-step nilsequences, which are functions of the form $F(g(n)\Gamma)$ where $G$ is a 2-step nilpotent Lie group. The proof for general $k$ uses a \emph{multiparameter} analysis on nilmanifolds, where the key step is to show that the nilmanifold must be abelian (hence a torus) in the case $k=2$, and more generally that the step of the nilpotent group is at most $k-1$.

\textbf{Random matrix universality.} With Van Vu, Tao proved the circular law (eigenvalues of random matrices with i.i.d. entries converge to the uniform distribution on the unit disk) and the Four Moment Theorem: eigenvalue statistics depend only on the first four moments of the entry distribution. This confirmed the Wigner--Dyson--Mehta universality conjecture for a broad class of random matrices. The proof of the circular law uses a \emph{local law} for the resolvent $(zI - X)^{-1}$ combined with a \emph{hopping argument} to move from the center of the spectrum to the edge.

The \emph{Four Moment Theorem} states that if $X$ and $Y$ are $n \times n$ random matrices with i.i.d. entries satisfying $\mathbb{E}[X_{ij}] = 0$, $\mathbb{E}[|X_{ij}|^2] = 1$, $\mathbb{E}[X_{ij}^3] = \beta_1$, and $\mathbb{E}[|X_{ij}|^4] = 3 + \beta_2$, then the joint distribution of the $k$ eigenvalues of $X$ closest to the origin converges to the same limit as for $Y$, provided $\beta_1$ and $\beta_2$ are the same for both matrices. The proof uses a \emph{moment matching} technique combined with a \emph{local Green's function estimate} for the resolvent $(zI - X)^{-1}$. The key insight is that the eigenvalue statistics are determined by the Green's function, and the Green's function depends on the entry distribution only through the first four moments.

The \emph{circular law} states that the empirical spectral measure of an $n \times n$ random matrix $X$ with i.i.d. entries of mean zero and unit variance converges weakly to the uniform distribution on the unit disk as $n \to \infty$. Tao and Vu proved this in full generality, settling a conjecture that had been open since Wigner's original work in the 1930s. The proof uses a \emph{local law} for the resolvent, a \emph{hopping argument} to move from the center of the spectrum to the edge, and a \emph{least squares argument} to control the error terms.

The Erd\H{o}s discrepancy problem. In 2015, Tao proved that for any sequence $x_1, x_2, \ldots$ with $x_i \in \{-1, 1\}$, the partial sums $\sum_{i=1}^n x_{id}$ are unbounded as $n \to \infty$ for any $d$, using Fourier analysis and the theory of completely multiplicative functions. The proof uses the \emph{entropy decrement} argument: at each step, the uncertainty (entropy) of the sequence increases, eventually forcing unbounded partial sums. The key insight is that for any completely multiplicative function $f: \mathbb{N} \to \{-1, 1\}$, the mean value $\mathbb{E}_{n \le N} f(n)$ tends to zero as $N \to \infty$, which forces the partial sums $\sum_{n \le N} f(n)$ to be unbounded.

The \emph{entropy decrement method} works as follows. Let $f: \mathbb{N} \to \{-1, 1\}$ be a completely multiplicative function. For each $k$, define the shifted function $f_k(n) = f(kn)$. The key observation is that the entropy of the joint distribution of $f(1), \ldots, f(N)$ increases by at least $\delta > 0$ at each step, where $\delta$ depends on the mean value of $f$. Since the entropy is bounded above by $\log 2$, the process must terminate after finitely many steps, which forces the mean value to zero and hence the partial sums to be unbounded.

\textbf{Energy-critical NLS.} With Colliander, Keel, Staffilani, and Takaoka, Tao proved global well-posedness and scattering for the energy-critical defocusing nonlinear Schr\"odinger equation $i\partial_t u + \Delta u = |u|^4 u$ in $\mathbb{R}^3$, using concentration compactness and Morawetz estimates. The key insight is a \emph{multilinear profile decomposition} that reduces the scattering problem to a study of non-interacting wave packets. The energy-critical NLS has the scaling symmetry $u(t,x) \mapsto \lambda u(\lambda^2 t, \lambda x)$, and the critical space is $H^1(\mathbb{R}^3)$. The defocusing condition ($|u|^4 u$ has a positive sign) ensures that the energy is a conserved quantity, and the proof of scattering uses the fact that the energy controls the $H^1$ norm globally.

The \emph{Strichartz estimates} for the Schr\"odinger equation, which Tao used extensively, provide $L^p_t L^q_x$ bounds for solutions of the free Schr\"odinger equation. These estimates are sharp and play the same role for dispersive PDEs that Sobolev estimates play for elliptic PDEs. The combination of Strichartz estimates with Morawetz estimates (which control the interaction between different spatial regions) provides the key tool for proving scattering for nonlinear dispersive equations.

\section{Impact of This Discovery}

The Green--Tao theorem transformed additive number theory and demonstrated the power of the transference principle, which has been applied to many other sparse sets including squarefree numbers, Chen primes, and polynomial progressions. The inverse theorem for Gowers norms has become a standard tool with applications to the Hardy--Littlewood circle method, the Chowla conjecture, and the Sarnak conjecture on the M\"obius function.

The Four Moment Theorem and the proof of universality have had significant impact in random matrix theory, nuclear physics, quantum chaos, and data science. The polynomial method, developed with Guth and Katz, solved the Erd\H{o}s distance problem and the joints problem in incidence geometry. Tao's proof of the circular law---that the empirical spectral measure of a random matrix with i.i.d. entries converges to the uniform distribution on the unit disk---settled a conjecture open since the 1930s and demonstrated the power of his local Green's function approach combined with a hopping argument across the spectrum.

Tao's Polymath projects pioneered massively collaborative mathematics, demonstrating that large-scale online collaboration can solve hard problems. Polymath8 optimized Zhang's bound on gaps between primes to 246, and Polymath1 tackled the density Hales--Jewett theorem.

The \emph{density Hales--Jewett theorem} (DHJT) states that for any $\delta > 0$ and $k \ge 2$, there exists $N = N(\delta, k)$ such that any subset of $[k]^N$ of density $\delta$ contains a \emph{combinatorial line}. The DHJT implies Szemer\'edi's theorem (by taking $k=2$) and is itself a deep combinatorial result. Tao and Gowers, independently and through Polymath1, gave new proofs using the transference principle and the inverse theorem for Gowers norms.

Compressed sensing, developed with Cand\`es, has had applications in medical imaging (faster MRI), digital cameras, astronomy, and radar through the restricted isometry property and $\ell_1$-minimization. The key result is that an $s$-sparse signal in $\mathbb{R}^n$ can be recovered from only $O(s \log(n/s))$ random measurements via $\ell_1$-minimization, which is far below the Nyquist rate. The \emph{restricted isometry property} (RIP) states that a measurement matrix $\Phi \in \mathbb{R}^{m \times n}$ satisfies $(1+\delta_s)$-RIP if for all $s$-sparse vectors $x$, $(1-\delta_s)\|x\|_2^2 \le \|\Phi x\|_2^2 \le (1+\delta_s)\|x\|_2^2$. Random Gaussian and sub-Gaussian matrices satisfy RIP with high probability when $m \ge C s \log(n/s)$, and this bound is essentially optimal. The compressed sensing framework has been extended to low-rank matrix recovery (using nuclear norm minimization), to structured sparsity, and to nonlinear inverse problems (phase retrieval, dictionary learning).

The \emph{polynomial method}, developed with Guth, provides a powerful tool for incidence geometry and combinatorial geometry. The key insight is that algebraic varieties of controlled degree have bounded intersection complexity, and this complexity can be used to bound the number of incidences between points and lines (or other geometric objects). The polynomial method solved the Erd\H{o}s distance problem (showing that $n$ points in $\mathbb{R}^2$ determine at least $n^{2/2-\varepsilon}$ distinct distances) and the joints problem (bounding the number of joints formed by lines in $\mathbb{R}^3$).

\section{Modern Views and Improvements}

The transference principle continues to be extended to new settings. Tao's blog, one of the most widely read mathematical blogs, has become an important resource for the mathematical community and a model for open-access mathematics.

The \emph{Chowla conjecture}---that the M\"obius function values $\mu(n_1), \ldots, \mu(n_k)$ are asymptotically independent for distinct $n_1, \ldots, n_k$---remains open, but Tao and Helfgott proved that the Sarnak conjecture (that $\sum_{n \le N} \mu(n) f(n \theta) \to 0$ for any deterministic sequence $f$) follows from the Chowla conjecture. This connects multiplicative number theory to ergodic theory in a fundamental way. The key insight is that the Chowla conjecture implies that the M\"obius function is ``orthogonal'' to any deterministic sequence, which is exactly the statement of the Sarnak conjecture.

The \emph{Sarnak conjecture} predicts that the M\"obius function is orthogonal to any bounded sequence $f(n)$ that arises from a topological dynamical system of zero entropy. Tao's work connects this conjecture to the Chowla conjecture by showing that the Chowla conjecture implies the Sarnak conjecture for all sequences of zero entropy. The proof uses the \emph{entropy decrement argument}: if the M\"obius function were correlated with a deterministic sequence, then the entropy of the sequence would increase, contradicting the zero-entropy assumption. This argument is a variant of the entropy decrement method used in the Erd\H{o}s discrepancy problem.

The energy-critical wave equation $\partial_t^2 u - \Delta u + u^5 = 0$ and the Navier--Stokes equations remain active areas. Tao's work on the Navier--Stokes equations suggests a possible route to proving finite-time blowup, though the problem remains one of the Clay Millennium Prize problems.

At 51 as of 2026, Tao continues to publish at an extraordinary rate, exploring AI-assisted mathematics, the solvability complexity index, and ongoing work on analytic number theory and the Riemann zeta function. Recent work on the \emph{Kakeya needle problem} and the \emph{Schr\"{o}dinger maximal function} builds on his earlier harmonic analysis contributions. Tao's work on the \emph{Navier--Stokes equations} has explored possible scenarios for finite-time blowup, though the Millennium Prize problem remains open.

The \emph{solvability complexity index} is a measure of the computational difficulty of a mathematical problem. Tao, with others, developed a framework for measuring the complexity of solving differential equations, integral equations, and other continuous problems using Turing machines. The key result is that the solvability complexity index of the Navier--Stokes equations in 3D is at least 2, meaning that any algorithm for deciding regularity must make at least two ``queries'' to the solution. This framework provides a rigorous way to measure the ``hardness'' of mathematical problems and connects analysis to theoretical computer science.

The \emph{Navier--Stokes equations} in $\mathbb{R}^3$ are the fundamental equations of fluid mechanics:
\[
\partial_t u + (u \cdot \nabla) u = \nu \Delta u - \nabla p, \qquad \nabla \cdot u = 0,
\]
where $u(t,x)$ is the velocity field, $p$ is the pressure, and $\nu > 0$ is the viscosity. The Millennium Prize problem asks whether solutions with smooth initial data remain smooth for all time, or whether finite-time singularities can develop. Tao's contribution is the construction of \emph{averaged Navier--Stokes flows} that concentrate energy on a fractal set in spacetime. While these flows are not exact solutions of the Navier--Stokes equations, they satisfy an averaged version and provide evidence that the regularity conditions may fail. The construction uses a \emph{convex integration} scheme combined with a \emph{microlocal analysis} of the energy cascade.

The \emph{convex integration} method, originally developed by Nash for isometric embeddings and refined by De Lellis and Sz\'ekelyhidi for the Euler equations, produces weak solutions of the Euler equations from any smooth ansatz by iteratively adding highly oscillatory perturbations. Tao adapted this method to the averaged Navier--Stokes equations, showing that the energy can be concentrated on sets of arbitrary Hausdorff dimension less than 5/3. The key insight is that the nonlinear term $(u \cdot \nabla) u$ can be used to transfer energy from large scales to small scales, and that the dissipation term $\nu \Delta u$ can be made to act only on a small set in spacetime.

The \emph{polynomial method} in combinatorics, developed with Guth, provides a powerful tool for incidence geometry. The key insight is that algebraic varieties of controlled degree have bounded intersection complexity, and this complexity can be used to bound the number of incidences between points and lines (or other geometric objects). The polynomial method solved the Erd\H{o}s distance problem (showing that $n$ points in $\mathbb{R}^2$ determine at least $n^{2-\varepsilon}$ distinct distances) and the joints problem (bounding the number of joints formed by lines in $\mathbb{R}^3$). The polynomial partitioning technique, which partitions $\mathbb{R}^n$ into cells using the zero set of a polynomial, is the key tool. The degree of the polynomial is chosen to balance the number of points in each cell, and the induction on dimension reduces the problem to a study of incidences on algebraic varieties.

The \emph{Schr\"{o}dinger maximal function} measures the local behavior of solutions to the free Schr\"{o}dinger equation $i\partial_t u + \Delta u = 0$. The sharp $L^2$ regularity theorem of Tao, Du, and Zhang determines the largest exponent $s$ such that the maximal function $\sup_{t > 0} |u(t,x)|$ is bounded in $L^2(\mathbb{R}^d)$ when $u_0 \in H^s(\mathbb{R}^d)$. This result connects harmonic analysis to dispersive PDE and has applications to the pointwise convergence problem for Schr\"{o}dinger equations.

Tao's work on the \emph{Navier--Stokes equations} includes the construction of fluid flows that concentrate energy on a fractal set in spacetime, providing evidence that the Millennium Prize problem's regularity conditions may fail. While this does not constitute a proof of blowup, it shows that any proof of regularity must use deep structural properties of the equations beyond standard energy estimates. The construction uses a \emph{pseudo-conformal conservation law} combined with a \emph{concentration compactness} argument to create solutions that concentrate energy on sets of Hausdorff dimension less than 5/3.

The \emph{Polymath project} model of collaborative mathematics, initiated by Tao and Gowers, has produced several breakthroughs. Polymath1 solved the density Hales--Jewett theorem, and Polymath8 improved the best known bound on prime gaps to 246 (unconditionally) and to 6 (assuming the Elliott--Halberstam conjecture). These projects demonstrated that large-scale online collaboration can solve problems that resist individual effort, and they have influenced the sociology of mathematical research.

\begin{center}\rule{0.5\linewidth}{0.5pt}\end{center}

\section*{Summary}
\addcontentsline{toc}{section}{Summary}

Tao proved the Green--Tao theorem that the primes contain arbitrarily long arithmetic progressions, developed the transference principle for sparse sets, proved the inverse theorem for Gowers norms, and established universality for random matrices. His work spans harmonic analysis, PDEs, combinatorics, number theory, and data science.
